\section{연습문제}

\begin{exercise}
$\Q[X]/(X^2-3)$의 원소를 표준형으로 나타내고 $(2+\alpha)^{-1}$을 구하여라. 여기서 $\alpha^2=3$이다.
\end{exercise}

\begin{exercise}
$X^4-5X^2+6$의 $\Q$ 위 분해체와 그 차수를 구하여라.
\end{exercise}

\begin{exercise}
$X^4-2$의 분해체가 $\Q(\sqrt[4]{2},i)$임을 보이고 그 차수를 구하여라.
\end{exercise}

\begin{exercise}
$X^3-1$의 $\Q$ 위 분해체와 차수를 구하여라.
\end{exercise}

\begin{exercise}
$f\in K[X]$가 이미 $K[X]$에서 일차인수의 곱으로 분해되면 그 분해체가 $K$임을 증명하여라.
\end{exercise}

\begin{exercise}
$f,g\in K[X]$의 분해체를 각각 $L_f,L_g$라 하자. $fg$의 분해체가 두 체의 합성체 $L_fL_g$와 $K$-동형임을 설명하여라.
\end{exercise}

\begin{exercise}
$X^p-X\in\F_p[X]$의 분해체를 구하여라.
\end{exercise}

\begin{exercise}
유한체는 대수적으로 닫혀 있지 않음을 증명하여라.
\end{exercise}

\begin{exercise}
$L/K$가 $f\in K[X]$의 분해체이고 $K\subseteq E\subseteq L$이라 하자. $L$이 반드시 $E$ 위에서 $f$의 분해체인지 판정하고 이유를 설명하여라.
\end{exercise}

\begin{exercise}
차수 $3$인 다항식의 분해체 차수가 가질 수 있는 값이 $1,2,3,6$뿐임을 보이고 각각의 예를 찾아라.
\end{exercise}
